% =====================================================================================================
\section{Are the assumptions reasonable?}\label{sec:assumptions}

\subsection{The answer}\label{sec:a-answer}

Our theorems use the same kinds of assumptions as \citet{SongLiang2015}: a Gaussian linear model, exact sparsity, a
bound on the dimension, a law for the predictors, eigenvalue floors on small sets of columns, and nonzero coefficients
that are not too small. On the designs we study, the assumptions of our posterior theorems hold, strong factors
included. Where the two papers differ, the direction depends on the regime --- their model-size cap, whether predictors
are centred, and which of their theorems is the comparator --- and is given assumption by assumption in the table.

For the coded method, Section~\ref{sec:who-outside} says which designs are outside Result~3 and what still
holds for them; conclusions are compared with Song and Liang theorem by theorem in Section~\ref{sec:SL-concl}, not as a
whole.

\subsection{Our designs, and what ``inside'' means}\label{sec:a-setting}

\paragraph{The designs.} All our designs have i.i.d.\ Gaussian rows, noise variance $1$ and fixed nonzero coefficients.
They fall into three classes.
\begin{itemize}[leftmargin=1.5em,itemsep=1pt]
  \item \emph{Independent:} S1 of the 2,400-fit study, and IND of the $10^6$ plan (planned runs with $p$ up to $10^6$).
  \item \emph{Locally correlated:} S2 (AR(1), correlation $0.5$), S3 and S4 (blocks of five at $0.7$ and $0.85$ respectively), S6 (blocks of
    five at $0.4$--$0.9$), AR (correlation $0.6$, the primary design of the $10^6$ plan) and BLK (blocks of $20$ at
    $0.5$).
  \item \emph{Pervasive factor:} S5 (blocks of five at $0.7$ plus a global factor, $0.3$ between all columns); the
    strong-factor study (one factor at $0.3$, $0.5$ or $0.7$; two factors; and its own version of S5; all true effects of
    one sign), with $n$ from $50$ to $150$ and $p$ up to $10n$; and FAC (one factor at $0.3$), a
    contingent design of the $10^6$ plan, run at $p=10^4$ only.
\end{itemize}
The 2,400-fit study has $n=200$, $p\in\{500,1000\}$ and $14$ true predictors, one in each group of five; the $10^6$ plan
has $5$ true predictors.

\paragraph{What ``inside'' means.} A design is \emph{inside} an assumption when the design, continued as a sequence in
$n$ with its structure and the number of true predictors fixed, satisfies it. We judge the qualitative conditions (a
floor bounded below, nonzero coefficients that do not change with $n$, a leakage set of size $n^b$); the explicit
constants of the proofs are conservative and are not part of the judgement. Page numbers are journal pages of their paper; ``supp.\ p.'' is a page of its online supplement; $\bar r_n$ is
their model-size cap.


\paragraph{How the conclusions compare.}\phantomsection\label{sec:SL-concl} Against their Theorem~3 (p.~956), our selection theorems add specificity and
exact selection at a polynomial rate, while theirs gives sensitivity with an exponentially small error; neither implies
the other. Against their Theorem~S (supp.\ p.~13), the statements concern different posteriors, so neither implies the other; within one posterior,
mass on the true model with a rate is a stronger form of statement than a statement about the posterior mode. The
distribution theorem uses the same distance as their Theorems~1--2, with the same exponential form of the tail for
levels $M\le\sqrt n/\log n$ of the radius $\sqrt{M\log n/n}$; it reaches a radius of order $\rn$, which their conditions do not reach at polynomial $p$;
larger radii and exponential $p$ are outside ours. For the reported model, their
statements are limits in the number of iterations at fixed data, for every data set; ours hold at the coded budget
under design conditions; neither implies the other.

\begin{table}[!t]
\centering\footnotesize
\setlength{\tabcolsep}{3.5pt}
\renewcommand{\arraystretch}{1.0}\setlength{\abovecaptionskip}{4pt}
\caption{The assumptions, \SL{} (2015), and our designs, by class: independent, locally correlated, pervasive factor.
\yes{} inside, \no{} outside, \unc{} unclear, \na{} not applicable. A comparison in one word appears only where it holds
in every regime of both papers. ``Their sets'': the column sets of their eigenvalue conditions, growing with their model-size cap
$\bar r_n$. ``No counterpart'': they assume nothing of the kind; ``neither implies the other'': both assume something,
unordered.}\label{tab:a-summary}
\begin{tabular}{@{}>{\raggedright\arraybackslash}p{0.18\linewidth}>{\raggedright\arraybackslash}p{0.16\linewidth}>{\raggedright\arraybackslash}p{0.43\linewidth}ccc@{}}
\toprule
Assumption & In words & \SL: counterpart, and comparison by regime & Indep. & Local & Factor \\
\midrule
\multicolumn{6}{@{}l}{\textbf{All results}} \\
Gaussian linear model \lab{GL} & Gaussian noise, fixed variance & (1), p.~947; (2), p.~950. \emph{The same, in every regime.} & \yes & \yes & \yes \\
Dimension \lab{DIM} & $p\le n^A$ & Cor.~1, p.~954; Thm~S, supp.\ pp.~13, 16; their $P_n\succ n^\theta$, p.~950.
Upper range: ours covers less. Lower range: ours covers bounded $p$, theirs does not. & \yes & \yes & \yes \\
Fixed true model size \lab{SP} & $\sz$ fixed & (4), p.~950. \emph{The same} for their Thms 2, 3, S; their Thm~1 also
admits growing supports under (3). & \yes & \yes & \yes \\
\midrule
\multicolumn{6}{@{}l}{\textbf{Results 1 and 2 (the posterior)}} \\
Signal size & bounded; nonzero (fixed in Thms \ref{thm:post-random}, \ref{thm:reported}, \ref{thm:outputs}; beta-min in Thm~\ref{thm:post-fixed}) & If their $\bstar$ does not change with $n$: the same as Thms \ref{thm:post-random}, \ref{thm:reported}, \ref{thm:outputs}; Thms \ref{thm:post-fixed}, \ref{thm:dist} admit more. If it may
vary: ours asks more than their proof steps use (supp.\ pp.~6, 15). & \yes & \yes & \yes \\
Eigenvalue floor \lab{EF}, \lab{EFS} & fixed floor on small sets & Cond.~3, p.~955; (22)--(23), supp.\ p.~13. Where their
sets are no larger than ours, ours implies (22)--(23) for standardised columns and Cond.~3 for Gaussian predictors; where they are larger ($\bar r_n\succ\log n$), neither implies the
other. & \yes & \yes & \yes \\
Near-exact collinearity \lab{NC} & floor of order $1/n$ on $O(\log n)$ sets & None in their Thms~1--2. Against Cond.~3, p.~955, and (22), supp.\ p.~13: theirs gives ours, on the sets they cover, where
their sets are at least as large; neither implies the other where they are smaller ($\bar r_n\prec\log n$). & \yes & \yes & \yes \\
Law of the rows \lab{GR} & Gaussian or sub-Gaussian & $\abs{x_i}\le1$, p.~952; Cond.~1, p.~953. Centred, unit-variance
predictors: ours implies Cond.~1. Raw predictors, or against $\abs{x_i}\le1$: neither implies the other. & \yes & \yes & \yes \\
\midrule
\multicolumn{6}{@{}l}{\textbf{Result 3 (the reported model), in addition}} \\
Visibility \lab{V} & no true predictor cancelled by $\le15$ others & Cond.~5, p.~958. \emph{Ours implies theirs, in every
regime.} & \yes & \yes & \unc$^{a}$ \\
Polynomial leakage set, \ref{asm:reported}(a) & at most $n^b$ nulls correlated with the signal & \emph{No hypothesis of this kind}; informal
remark (a), p.~958. No ordering. & \yes & \yes & \no \\
Floor and local regressions, \ref{asm:reported}(b),(c) & each regression in the set uses few members & Floor: as the eigenvalue floor,
on other sets. Local regressions: \emph{no counterpart}. & \yes & \yes & \na \\
\midrule
\multicolumn{6}{@{}l}{\textbf{Result~4 (Theorem~\ref{thm:outputs}), in addition}} \\
No leakage \lab{NL} & null predictors nearly uncorrelated with true ones & \emph{No counterpart}, in every regime. & \yes
& \no & \no \\
\bottomrule
\end{tabular}
\par\smallskip\raggedright\footnotesize
$^{a}$ Inside for every design of the strong-factor study (true effects of aligned signs); unclear for S5 and FAC.
\end{table}

\subsection{The assumptions one by one}\label{sec:a-each}

Every design is inside each assumption of Sections~\ref{sec:a-all} and~\ref{sec:a-post}; ``Our designs'' lines
appear only in Sections~\ref{sec:a-rm} and~\ref{sec:a-pe}. Each paragraph opens with the comparison with \SL.

\subsubsection{Assumptions of all results}\label{sec:a-all}

\ahead{Gaussian linear model}{\lab{GL}}{The same as \SL{} in every regime (their model (1), p.~947, and density (2), p.~950).}
\ev{Their density is $f^*(y,x)=\phi(y;x^\top\bstar,\sigs)\,\nu_x(dx)$; ours is the same model with an intercept.}
\ro{Heavy-tailed or heteroscedastic errors; no theorem of ours or of theirs covers them.}

\ahead{Dimension $p\le n^A$}{\lab{DIM}}{\emph{Upper range:} ours covers a smaller range of $p$ than \SL: far fewer than their Corollary~1
($P_n<\exp(Cn^\alpha)$, p.~954), slightly fewer than their Theorem~S (supp.\ p.~13). \emph{Lower range:} ours covers
bounded $p$; theirs needs $P_n\succ n^\theta$ (p.~950).}
\ev{Their Theorem~S needs $\log(\sqrt n/r_n)\succ\sqrt{\log P_n}$ (supp.\ p.~13, condition~2; $r_n\le\bar r_n$ is a prior parameter), hence
$\log P_n\prec(\log n)^2/4$; their Corollary~S$_1$ states $P_n<n^{\kappa\log n}$ with $\kappa\in(0,\tfrac14)$ (supp.\
p.~16). Our largest exponent is the $10^6$ plan at
$n=200$: $A=\log10^6/\log200\approx2.6$.}
\ro{Numbers of predictors that grow faster than every power of $n$, such as $\log p\asymp n^\alpha$.}

\ahead{Fixed true model size}{\lab{SP}}{The same as \SL's strict sparsity (4), p.~950, used by their Theorems~2, 3 and~S;
their Theorem~1 also allows growing supports under their (3), p.~950, which ours does not.}
\ev{Growth of $\sz$ is tightly limited: on every design with $\sz\le\min(K-1,p/2)$, $K$ the prior's cap on model
size, $\pis{\Sz}\le\{1+(\sz+1)/(3\sqrt n)\}^{-1}$
(Remark~\ref{rem:sharp}(a)), so a posterior miss of order $\rn$ is possible only if $\sz=O(\sqrt{\log n})$.}
\ro{True models whose size grows with $n$.}

\subsubsection{Assumptions of Results 1 and 2}\label{sec:a-post}

\ahead{Signal size}{\lab{BC}, \lab{FSb}, \lab{FSt} and Assumption~\ref{asm:dist}(e) bound the coefficients above; \lab{BM} is
the beta-min}{\SL{} state no lower bound in their theorems. \emph{If their $\bstar$ does not change with $n$} (the natural
reading of p.~950): Theorems~\ref{thm:post-random}, \ref{thm:reported} and~\ref{thm:outputs} ask the same; Theorems~\ref{thm:post-fixed} and~\ref{thm:dist}, whose
coefficients may vary with $n$, admit more. \emph{If it may vary with $n$:} their proof steps need rates (supp.\
pp.~6, 15), and our fixed signals ask more than those steps use.}
\ev{Our explicit beta-min \lab{BM}, at the usual order $\sigs\rn$,
appears only in the fixed-design Theorem~\ref{thm:post-fixed}; it has the same form as their proof rates, with a fixed multiple of
$\log n$ where theirs has a diverging multiple of $\bar r_n\log P_n$ (supp.\ p.~15). Their Theorem~1 bounds the
coefficients above through (3) and (8) (pp.~950, 952). Without any lower bound, specificity can fail
(Remark~\ref{rem:no-betamin}).}
\ro{True coefficients below a constant times $\rn$ (Theorem~\ref{thm:post-fixed} only), and signals that grow without bound.}

\ahead{Eigenvalue floor on small sets}{\lab{EF}, \lab{EFS}, Assumption~\ref{asm:margin}(a) and the floor of Theorem~\ref{thm:dist-param}}{\emph{Where \SL's sets are no larger than ours:} our fixed floor
implies their (22)--(23) (supp.\ p.~13; standardised columns) and their Condition~3 (p.~955; Gaussian predictors), and we
also ask a floor on sets where they ask nothing. \emph{Where their sets are larger} ($\bar r_n\succ\log n$): the two constrain different sets; neither implies
the other.}
\ev{``No larger'' means $\abs t+\bar r_n\le2(A+1)\sz+2$ ($\abs t$ their true-model size) against Theorem~\ref{thm:post-fixed}, and $\bar r_n+1\le5(A+1)(\sz+1)\log n$
against Theorem~\ref{thm:post-random}. Condition~3 asks a \emph{vanishing} level $\delta_n\succ\epsilon_n^2$ ($\epsilon_n$ their contraction rate) on $\bar r_n+1$ columns that contain a true
predictor. We ask
a \emph{fixed} level on sets of $(2A+2)\sz+2$ predictors (fixed design) or $O(\log n)$ predictors (Gaussian rows). In
every regime both conditions follow from $\lmin(\Sig)\ge c$. No upper eigenvalue bound is used, so pervasive factors are allowed: the smallest eigenvalue of
$\Sig$ is $1$ (S1, IND), $1/3$ (S2), $0.3$ (S3, S5), $0.15$ (S4), $0.1$ (S6), $1-\rho$ (one factor at $\rho$), $0.40$ (two
factors), $0.25$ (AR), $0.5$ (BLK) and $0.7$ (FAC).}
\ro{Exact or asymptotic duplicates of a true predictor, and correlations tending to one with a true predictor.}

\ahead{No near-exact collinearity on sets of logarithmic size}{\lab{NC}, Assumption~\ref{asm:dist}(d)}{\emph{Against their Theorems~1--2:} a condition they do not have. \emph{Against
their Condition~3 (p.~955) and (22) (supp.\ p.~13), if their sets are at least as large as ours}
($\bar r_n\ge3(A+1)(\sz+1)\log n$): theirs gives ours on the sets they cover. \emph{If their sets are smaller}
($\bar r_n\prec\log n$): ours constrains sets theirs does not; neither implies the other.}
\ev{Their eigenvalue conditions (11) and (14) (p.~952) are on the prior covariance. With their own ridge choice, a ridge
$c_n\propto1/n$ (p.~953), no design condition is left, and their rate carries the diverging factor $\bar r_n$:
$\epsilon_n^2\succ\bar r_n\log P_n/n$ (their (10)). Our floor, $c_\psi/n$
on the correlation scale, is of the same order as that ridge, on $O(\log n)$ sets only, and gives the distribution
theorem its contraction radius $\sqrt{M\log n/n}$, for a fixed level $M$. Without it the distribution theorem still
holds, with an extra factor $\sqrt{\log(\praw\vee n)}$ in the radius, $\praw$ the number of raw predictors (Proposition~\ref{prop:dist}(ii)).}
\ro{Duplicated columns (null duplicates included) and exact linear dependence, as in one-hot or compositional coding.}
\tu{Inside after removing exact duplicates and exact linear dependences.}

\ahead{Law of the rows}{\lab{GR}, Assumption~\ref{asm:dist}(c)}{On centred, unit-variance predictors, Gaussian rows and
Assumption~\ref{asm:dist}(c) each imply \SL's Condition~1 (p.~953); on raw predictors, and against their $\abs{x_i}\le1$
(Theorem~1, p.~952), neither implies the other.}
\ev{For standardised predictors, Assumption~\ref{asm:dist}(c) implies their Condition~1 for every $r$, with $\delta=\sqrt{3/8}/L$
and $M=\log2$, $L$ being the sub-Gaussian constant of Assumption~\ref{asm:dist}(c) (a two-line computation, independently
checked). On raw predictors, Gaussian rows with mean $n$
satisfy our row law and violate Condition~1, which is a raw exponential-moment bound. Our generators draw centred
Gaussian rows with unit variances, which satisfy Condition~1 itself.}
\ro{Non-Gaussian rows for Theorems~\ref{thm:post-random}, \ref{thm:reported} and~\ref{thm:outputs}; heavy-tailed predictors for the distribution theorem.}
\tu{Bounded data such as genotypes are outside the Gaussian-row theorems as stated, and inside \SL's Condition~1; they are
covered by the fixed-design Theorem~\ref{thm:post-fixed} when the observed design meets its conditions.}

\subsubsection{Further assumptions of Result 3}\label{sec:a-rm}

\ahead{Visibility}{\lab{V}}{Ours implies \SL's Condition~5 (p.~958: $\abs{\operatorname{cov}(x_t,y)}>0$ for every true
predictor) in every regime: theirs is our condition with no conditioning and no margin.}
\ev{Visibility asks $\abs{\Cov(W_j,W^\top\tstar\mid W_U)}\ge a>0$ for every true predictor $j$, with $W$ a standardised row of the design and $U$ any set of at most $15$ other columns: the start law (the distribution the sampler starts from) must
see every true predictor, as their first stage must, and it conditions on up to $15$ columns because the start law works on groups of at most $16$ predictors. For bounded blocks with one true predictor each, and for AR
designs, this covariance is $\theta^*_j\Var(W_j\mid W_U)$ plus geometrically small terms, at
least $\abs{\theta^*_j}\lmin(\Sig)/2$. When the true effects have aligned signs and every conditional covariance
between columns outside $U$ is nonnegative, as in each design of the strong-factor study, it is at least
$\theta^*_j\lmin(\Sig)$.}
\ro{Masking: a true predictor whose signal some $15$ other predictors cancel.}
\ds{\textbf{Inside:} the independent and locally correlated designs, and every design of the strong-factor study.
\textbf{Unclear:} S5 (mixed signs; a search found a near-cancellation) and FAC (its signs are mixed, so the aligned-sign
bound does not apply; a search found a margin below it).}
\tu{A numerical check of the conditions' population quantities, not yet complete, found visibility failing for a
pervasive factor with effects of mixed signs and for a riboflavin-like design; such a user may be outside Result~3
and still inside Results~1 and~2.}

\ahead{Polynomial leakage set,}{Assumption~\ref{asm:reported}(a)}{\SL{} have no hypothesis of this kind; the nearest text is their
informal remark (a), p.~958, which is not a condition of any of their theorems. No ordering is claimed.}
\ev{They write: ``if there are $p_n$ predictors correlated with the true predictor $X_t$, $p_n\prec P_n$, then there are
at most $cp_n$ predictors to be selected'' (p.~958). Their results on their sampler's output carry no design condition of
this kind (supp.\ pp.~9--10, 12--13, 16--17). We state the requirement as a hypothesis with a rate: the null predictors
whose conditional covariance with a true predictor exceeds a level $\varpi_n=o(\rn)$ must lie in a set of at most $n^b$
columns, the leakage set $Q$, with $b<\tfrac12$.}
\ro{Pervasive factors with $p\gg\sqrt n$: with no conditioning, every null predictor then has a fixed covariance with every
true one, so the leakage set is all of $\set{p}$, and $\abs{Q}\le n^b$ would force $p\le n^b$.}
\ds{\textbf{Inside:} S1 (no leakage set beyond the true model); S2 and AR (a neighbourhood of size $C\log n$ around the true
predictors); S3, S4, S6 and BLK (the true predictors' blocks). \textbf{Outside:} S5, the whole strong-factor study, and
FAC.}

\ahead{Floor and local regressions inside the leakage set,}{Assumption~\ref{asm:reported}(b),(c)}{\emph{Floor on $Q\cup\{l\}$,
for every predictor $l$:} of the kind of \SL's Condition~3
(p.~955) and (22) (supp.\ p.~13), on other sets; it compares as the eigenvalue floor does. \emph{Local regressions:} no
counterpart (they serve the method's local search, which \SL{} do not analyse).}
\ev{Both follow from the global floor \lab{EFS} plus a locality property of $\Sig$: for block, banded and AR structures,
take as local set the block-mates or the nearest leakage-set neighbours. A fixed-design version needs nothing of the kind
when the leakage set has at most $(1-\delta)\log n$ extra predictors, for a fixed $\delta\in(0,1)$ (Theorem~\ref{thm:reported-fixed}, case
(d1)).}
\ro{Leakage sets collinear with another predictor, or inside which regressions depend on many members.}
\ds{As for the polynomial leakage set; not applicable where that fails.}

\subsubsection{Further assumption of Result 4}\label{sec:a-pe}

\ahead{No leakage}{\lab{NL}}{In no regime do \SL{} have a condition of this kind: none of their theorems concerns a
start law. Only our independent designs satisfy it.}
\ev{With no conditioning, \lab{NL} asks $\abs{\Cov(W_l,W_j)}\le C_Ln^{-1/2}$ for every null predictor $l$ and true
predictor $j$. In AR(1) with correlation $0.5$ the neighbour of a true predictor has covariance $0.5$; in a block design
every block-mate of a true predictor has covariance at least $0.4$. For such designs the start law concentrates on the
models between the true model and the leakage set, not on the true model alone. On AR and BLK the present proof route
cannot be extended as $n\to\infty$, because the start law's mass on the true model then decays faster than any power of
$n$; at the sizes of the $10^6$ plan this asymptotic decay says nothing, and we have no bound on this mass there in either
direction. For S3, S4 and S6, a theorem under a weaker, group-relative form of \lab{NL} is a candidate; it is not
proved.}
\ro{Any null predictor correlated with a true one at a fixed level.}
\ds{\textbf{Inside:} the independent designs only. \textbf{Outside:} all locally correlated and factor designs, including AR, the
primary design of the $10^6$ plan.}


\subsection{Who is outside Result 3, and what still holds}\label{sec:who-outside}

The assumptions of Results~1 and~2 hold on every design of our studies, strong common factors included; those of Result~3 hold on the independent
and locally correlated designs (S1--S4, S6, IND, AR, BLK) and not on the pervasive-factor designs (S5, the strong-factor
study, FAC): with $p\gg\sqrt n$, every null predictor then has a fixed covariance with every true one, so the leakage set
is all of $\set{p}$ (Section~\ref{sec:a-rm}). Visibility can also fail on its own when true effects of mixed signs cancel;
a numerical check of the conditions' population quantities, not yet complete, found this for a pervasive factor with
effects of mixed signs and for a riboflavin-like design. Theorem~\ref{thm:outputs} of Result~4 needs no leakage at all,
which only the independent designs satisfy. What still holds outside:
\begin{itemize}[leftmargin=1.5em,itemsep=1pt]
  \item on the block and AR designs: Results~1, 2 and~3, and the unbiasedness of the evidence estimate
    (Proposition~\ref{prop:unbiased}); no theorem yet gives the consistency of the reported inclusion probabilities,
    median model and evidence estimate there;
  \item on the pervasive-factor designs: Results~1 and~2 and Proposition~\ref{prop:unbiased}; certified screening
    (Theorem~\ref{thm:screen}), which has no leakage condition, applies when its own conditions hold.
\end{itemize}
